% !TEX root = ../main.tex
\chapter{Giới thiệu}
\label{Chapter1}

\section{Bối cảnh và động lực nghiên cứu}
\label{sec:boi-canh}

Trong thập kỷ qua, hệ thống tư vấn (recommender system) đã trở thành một thành phần quan
trọng của nhiều nền tảng số, đặc biệt trong thương mại điện tử, dịch vụ phát trực tuyến và
mạng xã hội. Nhiệm vụ cốt lõi của hệ thống tư vấn là ước lượng mức độ phù hợp giữa người
dùng và sản phẩm từ dữ liệu tương tác lịch sử, qua đó cung cấp các đề xuất được cá nhân
hóa~\cite{zhang2019deep}.

Lọc cộng tác (collaborative filtering -- CF) là một trong những nền tảng quan trọng của các
hệ thống tư vấn hiện đại. Trong đó, phân rã ma trận (matrix factorization -- MF)~\cite{mf2009}
biểu diễn người dùng và sản phẩm bằng các véc-tơ ẩn trong một không gian đặc trưng thấp
chiều, sao cho mức độ tương thích giữa hai biểu diễn phản ánh xu hướng ưa thích của người
dùng. Các phương pháp mở rộng như LightGCN~\cite{lightgcn2020} tiếp tục khai thác cấu
trúc đồ thị của dữ liệu tương tác và đã cho thấy hiệu quả trên nhiều bộ dữ liệu chuẩn. Tuy
nhiên, các phương pháp dựa chủ yếu vào tín hiệu cộng tác vẫn đối mặt với ba hạn chế chính.
Thứ nhất, dữ liệu tương tác thường rất thưa, khiến biểu diễn của các người dùng hoặc sản
phẩm có ít tương tác kém ổn định. Thứ hai, bài toán khởi động nguội (cold-start) gây khó
khăn cho việc xây dựng biểu diễn có chất lượng đối với người dùng hoặc sản phẩm mới. Thứ
ba, các biểu diễn cộng tác không trực tiếp khai thác thông tin ngữ nghĩa từ nội dung văn bản
mô tả sản phẩm hoặc sở thích người dùng.

Sự phát triển nhanh chóng của các mô hình ngôn ngữ lớn (large language model -- LLM) mở
ra một hướng tiếp cận bổ sung cho các hạn chế trên. Nhờ được tiền huấn luyện trên quy mô
dữ liệu văn bản lớn, LLM sở hữu khả năng biểu diễn ngữ nghĩa và suy luận phong phú. Việc
kết hợp LLM với hệ thống tư vấn hiện được nghiên cứu theo hai hướng chính~\cite{llmrec2023}.
Hướng thứ nhất, LLM-enhanced, sử dụng LLM như một bộ mã hóa để làm giàu biểu diễn người
dùng và sản phẩm, sau đó đưa các biểu diễn này vào mô hình gợi ý truyền thống. Hướng thứ
hai, LLMs-as-RS, biểu diễn trực tiếp bài toán gợi ý dưới dạng một tác vụ ngôn ngữ, trong đó
LLM vừa xử lý ngữ cảnh vừa sinh đầu ra gợi ý. Hướng LLMs-as-RS đặc biệt đáng chú ý vì có
khả năng tận dụng tri thức ngữ nghĩa có sẵn của LLM trong các tình huống cold-start và hỗ
trợ diễn giải kết quả bằng ngôn ngữ tự nhiên.

Luận văn này tập trung vào hướng LLMs-as-RS, với mục tiêu khai thác đồng thời thông tin
ngữ nghĩa từ LLM và tín hiệu cộng tác từ lịch sử tương tác người dùng--sản phẩm.

\section{Vấn đề đặt ra: khoảng cách ngữ nghĩa giữa tín hiệu cộng tác và LLM}
\label{sec:van-de}

Khi bài toán gợi ý được biểu diễn như một tác vụ ngôn ngữ, một câu hỏi quan trọng là làm
thế nào để đưa lịch sử hành vi của người dùng vào LLM một cách hiệu quả. Cách trực tiếp
nhất là chuyển lịch sử tương tác sang văn bản, chẳng hạn liệt kê các sản phẩm hoặc bộ phim
mà người dùng đã tương tác, rồi bổ sung phần này vào chỉ dẫn đầu vào. Tuy nhiên, đối với
người dùng có lịch sử dài, cách biểu diễn này nhanh chóng làm tăng số lượng token và có thể
vượt quá giới hạn ngữ cảnh của LLM, đồng thời làm tăng đáng kể chi phí tính toán~\cite{collm2023}.

Một hướng thay thế là nén thông tin hành vi thành biểu diễn véc-tơ. Thay vì đưa toàn bộ lịch
sử dưới dạng văn bản, mô hình cộng tác tổng hợp lịch sử tương tác thành một véc-tơ dày đặc,
sau đó chuyển véc-tơ này sang không gian embedding của LLM và đưa vào như một soft token.
Cách tiếp cận này cho phép truyền tín hiệu cộng tác vào LLM mà không làm tăng đáng kể độ
dài ngữ cảnh.

CoLLM~\cite{collm2023} là một công trình tiêu biểu theo hướng này. Mô hình sử dụng một
MLP để ánh xạ biểu diễn cộng tác từ không gian của mô hình truyền thống sang không gian
embedding của LLM, sau đó chèn kết quả vào prompt dưới dạng soft token. Cơ chế này đơn
giản và hiệu quả, nhưng phép ánh xạ được áp dụng độc lập trên từng véc-tơ cộng tác. Mỗi
biểu diễn đầu vào được chuyển thành một soft token thông qua cùng một hàm ánh xạ, nên
mô-đun này không trực tiếp mô hình hóa quan hệ giữa nhiều nguồn tín hiệu cộng tác trước khi
chúng được đưa vào LLM, cũng như không có cơ chế chọn lọc thích nghi theo ngữ cảnh.

Thách thức cốt lõi nằm ở sự khác biệt giữa hai không gian biểu diễn. Biểu diễn cộng tác được
học từ cấu trúc tương tác người dùng--sản phẩm, trong khi không gian của LLM được hình
thành chủ yếu từ dữ liệu ngôn ngữ. Do nguồn gốc huấn luyện khác nhau, việc chuyển đổi giữa
hai không gian không chỉ là bài toán thay đổi số chiều, mà còn là bài toán căn chỉnh biểu
diễn. Từ đó, luận văn đặt ra câu hỏi: liệu một mô-đun căn chỉnh có khả năng chọn lọc và tổng
hợp tín hiệu cộng tác theo từng mẫu có thể cải thiện quá trình tích hợp tín hiệu cộng tác vào
LLM so với phép ánh xạ pointwise bằng MLP hay không?

\section{Khoảng trống nghiên cứu}
\label{sec:khoang-trong}

Q-Former, được giới thiệu trong BLIP-2~\cite{blip2_2023}, là một kiến trúc trung gian được
thiết kế để kết nối đặc trưng ảnh với mô hình ngôn ngữ. Thành phần này sử dụng một tập truy
vấn học được (learned query tokens) kết hợp với cross-attention để khai thác có chọn lọc thông
tin từ đặc trưng đầu vào. Thay vì chuyển toàn bộ đầu vào thông qua một phép chiếu trực tiếp,
các truy vấn học cách thu nhận và tổng hợp những phần thông tin phù hợp cho nhiệm vụ phía
sau.

Về nguyên tắc, cơ chế này cũng phù hợp với bài toán tích hợp biểu diễn cộng tác vào LLM.
Biểu diễn cộng tác có thể thay thế đặc trưng ảnh ở phía đầu vào, trong khi các truy vấn của
Q-Former đóng vai trò tổng hợp và biến đổi tín hiệu trước khi chuyển sang không gian của
LLM. ILM~\cite{ilm2024} đã khai thác ý tưởng này trong bài toán gợi ý, cho thấy Q-Former
có thể được sử dụng như một mô-đun trung gian giữa thông tin cộng tác và mô hình ngôn ngữ.

Tuy nhiên, vẫn còn hai khoảng trống liên quan trực tiếp đến mục tiêu của luận văn. Thứ nhất,
ILM không được đánh giá theo bài toán dự đoán CTR với AUC và uAUC như CoLLM và các
công trình cùng hướng, nên chưa tạo ra cơ sở so sánh trực tiếp giữa cơ chế Q-Former và phép
ánh xạ MLP trong cùng thiết lập đánh giá. Thứ hai, mô-đun Q-Former trong ILM chưa được
khảo sát cùng các kỹ thuật huấn luyện nhằm cải thiện tới chất lượng phân biệt giữa người dùng
và chất lượng xếp hạng trong từng người dùng (AUC/uAUC), cũng như lựa chọn biểu diễn nhiều soft token.

Từ các khoảng trống trên, luận văn xây dựng CoQLLM, trong đó Q-Former được sử dụng làm
mô-đun căn chỉnh giữa biểu diễn cộng tác và LLM. Mô hình được đánh giá trên bài toán dự
đoán CTR với AUC và uAUC, đồng thời được bổ sung cơ chế điều kiện hóa truy vấn theo
người dùng và biểu diễn cộng tác đa token.

\section{Mục tiêu và đóng góp}
\label{sec:dong-gop}

Luận văn hướng đến hai mục tiêu chính.

\textbf{Mục tiêu 1 -- Xây dựng mô-đun căn chỉnh Q-Former cho tín hiệu cộng tác.} Luận văn
thiết kế và hiện thực CoQLLM với ba thành phần chính: mô hình MF đóng băng, Q-Former làm
mô-đun căn chỉnh, và LLM Vicuna-7B-v1.5~\cite{vicuna2023} kết hợp LoRA~\cite{lora2021}.
Tín hiệu cộng tác được xử lý bởi Q-Former trước khi được đưa vào LLM dưới dạng soft token.
Trên kiến trúc cơ sở này, luận văn khảo sát điều kiện hóa truy vấn theo người dùng và biểu
diễn đa token.

\textbf{Mục tiêu 2 -- Đánh giá Q-Former cho bài toán CTR.} Luận văn đánh giá CoQLLM
trên MovieLens-1M và Amazon-Book bằng AUC và uAUC, đồng thời đối chiếu với các baseline
như CoLLM, BinLLM và CoRA trong phạm vi số liệu tương thích.

Từ hai mục tiêu trên, các đóng góp chính của luận văn gồm:

\begin{enumerate}
    \item \textbf{Thích nghi Q-Former cho bài toán gợi ý dựa trên tín hiệu cộng tác.} Đặc
    trưng cộng tác từ MF được sử dụng thay cho đặc trưng ảnh trong cơ chế cross-attention của
    Q-Former và được chuyển thành soft token đầu vào cho LLM. Trên nền tảng này, luận văn bổ
    sung điều kiện hóa truy vấn theo người dùng và khảo sát biểu diễn đa token.

    \item \textbf{Đánh giá có hệ thống theo AUC và uAUC.} CoQLLM được đánh giá trên
    MovieLens-1M và Amazon-Book theo thiết lập dự đoán nhị phân, cho phép đối chiếu trực tiếp
    hơn với dòng phương pháp CoLLM. Bên cạnh kết quả tổng thể, luận văn thực hiện phân tích
    thành phần và phân tích warm/cold để làm rõ vai trò của từng cơ chế và giới hạn của biểu
    diễn cộng tác.
\end{enumerate}

\section{Bố cục luận văn}
\label{sec:bo-cuc}

Phần còn lại của luận văn được tổ chức như sau. Chương~\ref{Chapter2} trình bày các công
trình liên quan, gồm lọc cộng tác truyền thống, các hướng khai thác LLM cho hệ thống tư vấn,
các phương pháp tích hợp biểu diễn cộng tác vào LLM và cơ sở kiến trúc Q-Former.
Chương~\ref{Chapter3} mô tả phương pháp đề xuất CoQLLM, bao gồm kiến trúc tổng thể,
mô-đun MF, Q-Former, cơ chế soft token, các kỹ thuật huấn luyện bổ sung và quy trình huấn
luyện ba giai đoạn. Chương~\ref{Chapter4} trình bày thiết lập thực nghiệm, kết quả so sánh,
phân tích thành phần, phân tích warm/cold và thảo luận. Chương~\ref{Chapter5} tổng kết các
đóng góp, nêu các hạn chế và đề xuất hướng nghiên cứu tiếp theo.
